%*************************************************
% In this file the first few pages are typeset.
% Make the changes accordingly
%*************************************************

% شماره صفحات با حروف
\pagenumbering{adadi}

%***************************
% 1st page: Blank
%***************************
\thispagestyle{empty}
\mbox{}
\pagebreak

%***************************
% 2nd page: Besmelah
%***************************
\thispagestyle{empty}
\begin{center}
	~\vfill
	\includegraphics[scale=1]{besm.jpg}
	~\vfill
\end{center}
\pagebreak

%***************************
% 3rd page: Title
%***************************
\thispagestyle{empty}
%\pagenumbering{gobble}
\newgeometry{left=3cm,right=3cm,top=2cm}
\begin{center}
\includegraphics[height=3cm]{iut_logo.png}
\vspace{0.4cm}

\textbf{دانشگاه صنعتی اصفهان}\\
\vspace{0.4cm}

{\large

	دانشکده مهندسی برق و کامپیوتر
}
\vspace{3.5cm}

{\Large
	\textbf{بهبود کیفیت و سرعت یادگیری در سیستم‌های چندعامله با استفاده از
	\\
	معیار جدید خبرگی و انتگرال فازی}\\
}
\vspace{2.5cm}

{\Large
	پایان‌نامه کارشناسی ارشد مهندسی کامپیوتر -- هوش‌مصنوعی و رباتیک\\
}
\vspace{1cm}

{\large
	\textbf{داریوش حسن‌پور آده}\\
}
\vspace{2.5cm}

{\large
	استاد راهنما\\
}
\vspace{0.5cm}

{\large
	\textbf{دکتر مازیار پالهنگ}\\
}
\vfill

\textbf{پاییز ۱۳۹۵}

\end{center}
\restoregeometry
\pagebreak

%***************************
% 4th page: Signatures
%***************************
\thispagestyle{empty}
\newgeometry{left=3cm,right=3cm,top=2cm}
\begin{center}
\includegraphics[height=3cm]{iut_logo.png}
\vspace{0.4cm}

\textbf{دانشگاه صنعتی اصفهان}\\
\vspace{0.4cm}

{\large
	دانشکده مهندسی برق و کامپیوتر
}
\vspace{1.8cm}

\vfill

{\Large
	پایان‌نامه کارشناسی ارشد رشته مهندسی کامپیوتر -- هوش‌مصنوعی و رباتیک آقای داریوش حسن‌پور آده\\
	\vspace{.3cm}
	تحت عنوان\\
}


\end{center}

\begin{center}
{\large
	\textbf{بهبود کیفیت و سرعت یادگیری در سیستم‌های چندعامله با استفاده از\\
	معیار جدید خبرگی و انتگرال فازی}
}
\end{center}
\vspace*{1.8cm}
در تاریخ ۱۳۹۵/۱۰/۲۰ توسط کمیته تخصصی زیر مورد بررسی و تصویب نهایی قرار گرفت:\\

{\normalsize
	
	\begin{tabular}{rr}
	\vspace*{.8cm}
	1- استاد راهنمای پایان‌نامه  & \hspace{2cm} دکتر مازیار پالهنگ \\
	\vspace{.8cm}
	3-استاد داور &\hspace{2cm} دکتر عبدالرضا میرزایی \\
	\vspace{.8cm}
	۴-استاد داور &\hspace{2cm} دکتر محمد حسین منشئی \\
	\vspace{.8cm}
	سرپرست تحصیلات تکمیلی دانشکده &\hspace{2cm} دکتر محمد رضا تابان \\
	\end{tabular}
}
\restoregeometry
\pagebreak

%***************************
% 5th page: Acknowledgment
%***************************
\thispagestyle{empty}
\newgeometry{left=3cm,right=4cm,top=4cm}
\vspace*{1.5cm}

{\large
	\textbf{تشکر و قدردانی}\\

	
پروردگار منّان را سپاسگزارم ......

}
\restoregeometry
\pagebreak

%***************************
% 6th page: Rights
%***************************
\thispagestyle{empty}
\newgeometry{left=6cm,right=6cm}

\begin{spacing}{3}
\leavevmode
\vfill
\parbox{8 cm}{

\textbf{\Large کلیه حقوق مادی مترتب بر نتایج مطالعات، ابتکارات و نوآوری‌های ناشی از تحقیق موضوع این پایان‌نامه متعلق به دانشگاه صنعتی اصفهان است.}

}
\vfill
\end{spacing}
\restoregeometry
\pagebreak

%***************************
% 7th page: Dedication
%***************************
\thispagestyle{empty}
\newgeometry{left=4.5cm,right=4.5cm,top=0cm,bottom=0cm}

%{\LARGE
%\centering
%\textbf{تقدیم به \\ پدر و مادر عزیزم }
%
%}
\bgroup\vspace*{\fill}
\renewcommand{\\}{\vspace{-.5em}\newline\relax}
\newcommand{\تب}[1][.5]{\hspace*{#1cm}}
\noindent
دلتنگی‌های آدمی را باد ترانه‌ای می‌خواند\\
\تب روياهايش را آسمان پر ستاره ناديده می‌گيرد\\
\تب[1] و هر دانه‌ی برفی به اشكی نريخته می‌ماند.\\
\تب[1.5] سكوت سرشار از سخنان ناگفته ا‌ست؛\\
\تب[2] از حركات نا‌كرده،\\
\تب[2.5] اعتراف به عشق‌های نهان،\\
\تب[3] و شگفتی‌های به زبان نيامده،\\
\تب[3.5] دراين سكوت حقيقت ما نهفته است؛\\
\تب[6] حقيقت تو و من.\\
\noindent
برای تو و خويش\\
\تب چشمانی آرزو می‌كنم،\\
\تب[1] كه چراغ‌ها و نشانه‌ها را در ظلمات‌مان ببيند.\\
\تب[1.5] گوشی،\\\تب[2] كه صداها و شناسه‌ها را در بيهوشی‌مان بشنود.\\
برای تو و خويش،\\
\تب[1.5] روحی،\\\تب[2] كه اين‌همه را در خود گيرد و بپذيرد.\\
\تب[1.5] و زبانی\\\تب[2] كه در صداقت خود ما را از خاموشی خويش بيرون كشد،\\
\تب[3] و بگذارد از آن‌چيزها كه در بندمان كشيده‌ است، سخن بگوييم.\\
\\
پنجه درافکنده‌ایم با دست‌هایمان\\
\تب[1] به جای رها شدن\\
سنگین سنگین بر دوش می‌کشیم\\
\تب[1] بار دیگران را\\
\تب[2] به جای همراهی کردن‌شان!\\
\تب[1]عشق ما نیازمند رهایی است نه تصاحب\\
 در راه خویش ایثار باید نه انجام وظیفه...\\
\\
\تب[3] بی‌اعتمادی دری است\\
 \تب[2.4] خودستایی، چفت و بست غرور است\\
\تب[1.8] و تهی دستی، دیوار است و لولاست\\
\تب[1.2] زندانی را که در آن محبوس رآی خویش‌ایم\\
\تب[.6] دلتنگی‌مان را برای آزادی و دلخواه دیگران بودن\\
از رخنه‌هایش تنفس می‌کنیم...\\
\begin{flushleft}\textit{
- مارگوت بیگل
}\end{flushleft}
\vspace*{\fill}\egroup
\pagebreak

%***************************
% 8th page: Table of contents
%***************************

\titleformat{\chapter}[display]
	{\normalfont\LARGE\bfseries\centering}{\chaptertitlename ~ \tartibi{chapter}}{20pt}{\LARGE}
\newgeometry{left=2.5cm,right=3cm,top=3cm,bottom=2cm,includehead=false,headsep=1cm,footnotesep=.5cm}
\baselineskip=.7cm

\addtocontents{toc}{\textbf{\underline{عنوان}}}
\addtocontents{toc}{\hfill\textbf{\underline{صفحه}}\par}
\addcontentsline{toc}{section}{فهرست مطالب}
\tableofcontents
\pagebreak

\addcontentsline{toc}{section}{فهرست تصاویر}
\listoffigures
\pagebreak

\addcontentsline{toc}{section}{فهرست جداول}
\listoftables
\pagebreak

% change the font and margins of a chapter title
\titlespacing*{\chapter}{0pt}{1.5cm}{6cm}
\titleformat{\chapter}[display]
	{\normalfont\LARGE\bfseries\raggedright}{\chaptertitlename ~ \tartibi{chapter}}{20pt}{\LARGE}

% No page numbers on the first page of a chapter
\assignpagestyle{\chapter}{empty}
